\lead{MoE and looping at matched data.} We retrained the two A5 architectures on the DNA pool: a 4-expert learned top-1 MoE with quantile balancing (about four times the FFN parameters of d8, the same active compute) and a looped model (prelude, a shared block applied 6 times with RMS normalization and recall, coda; fewer parameters than d8). Table~\ref{tab:p1ours} lists the results; the LR .01 and .001 trials sit at the end of the queue and are marked cut where unfinished.
\begin{itemize}
\item Neither beats dense. At LR .003 the MoE reaches 1.2656 (+0.003 over d8) and the loop 1.2693 (+0.007). The MoE tracks our dense runs through pass 24 (1.2753 against 1.2754 for dropout .05) and then gains less during the final LR decay, consistent with the extra capacity fitting the repeated 30k chunks rather than new structure.
\item This reverses the A5 ordering, where the MoE won clearly on fresh text. On 30k repeated DNA chunks the supplied suite is already flat above d6, so more parameters (MoE) buy nothing and fewer (loop) cost a little; the data, not the architecture, sets the floor.
\item Prediction we make for Part (d): any architecture change that only adds capacity will land within $\pm0.005$ of d8 here; only changes that act as regularizers (dropout, ensembles) move the number.
\end{itemize}
